\section{Discussion}
\label{sec:discussion}

\paragraph{Two mechanisms, two places to look.}
The representation mechanism shows up before any rollout: a position read-out on \emph{real}
encodings fails, decoded real frames lose the cube, and no change to the predictor or planner can
bring it back. The planning mechanism shows up only after the state has been kept: the gate latent
supports control from real transitions but its imagination does not. A practitioner who sees a
failed imagination policy should first check which of the two they are looking at, because the
remedies are disjoint. Detaching the target, adding decoders, or changing what the encoder is asked
to keep addresses the first; distributional or value-aware prediction addresses the second.

\paragraph{What the flow reference does and does not show.}
The flow model differs from the JEPA in supervision stage, not only in sampling versus point
prediction, so its advantage where it has one is a recipe-level result. The ordering claim survives
this caveat: a sampled predictor cannot restore a state that the encoder never kept, and where the
JEPA representation held the state, the remaining gap is the kind that sampling addresses. We do
not claim that generative modeling is necessary; for a quadratic goal cost a mean and a variance
suffice, and value-equivalent models can carry the needed expectations directly.

\paragraph{Honest boundaries.}
The empirical per-frame-noise-plus-wind failure has high prediction loss and is not an instance of
the zero-loss nuisance optimum of Proposition~\ref{prop:representation}; the proposition explains
why non-collapse cannot protect the state, not which finite-budget model drops it. The gate JEPA
encoder used for the behavioral comparison is model seed 1, chosen because it was the only seed on
which decoded controls had succeeded in development, and every flow gate policy shares one dynamics
checkpoint. Decoded-latent results are statements about information recoverable by a supervised
decoder, not about the unsupervised model alone. The worlds are small by design; Craftax shows the
consequence at the scale of a real game but under a fixed ten-minute budget.

\paragraph{Outlook.}
General world models need a principle for choosing what to keep, and predictability is not it.
The results here argue for objectives that price the controlled state's innovations as
information to model rather than as loss to avoid, and for evaluating a learned world model at
three separate points: what its representation retains, what its predictions preserve, and what a
policy trained inside it can do.
